\section*{अध्याय \ref{ch:b3:environmental-toxicology} --- पर्यावरण रसायन और विषविज्ञान}

\begin{solution}{exo:b3:environmental-toxicology:1}
किसी समनाभिकीय द्विपरमाणुक का एकमात्र कंपन द्विध्रुव आघूर्ण को शून्य बनाए रखता है,
और आर्गन में कोई कंपन नहीं होता: इनमें से कोई अवरक्त विकिरण का अवशोषण नहीं कर सकता। $\ce{CO2}$ अपनी
बंकन विधा (और अपने असममित तनन) से अवशोषण करता है, जो एक दोलायमान
द्विध्रुव उत्पन्न करते हैं; बंकन पृथ्वी के उत्सर्जन के ठीक मध्य में पड़ता है।
\end{solution}

\begin{solution}{exo:b3:environmental-toxicology:2}
$10^{0.1} = 1.26$: $[\ce{H+}]$ 26\,\% बढ़ता है।
\end{solution}

\begin{solution}{exo:b3:environmental-toxicology:3}
$(2.0 + 0.50)\ \unit{mmol/L} \times \qty{100.1}{mg/mmol} = \qty{250}{mg/L}$ $\ce{CaCO3}$।
\end{solution}

\begin{solution}{exo:b3:environmental-toxicology:4}
N: $0.60/14.0 = \qty{0.043}{mmol/L}$; P: $0.010/31.0 = \qty{3.2e-4}{mmol/L}$; N\,:\,P $= 133$, 16 से बहुत ऊपर: फ़ॉस्फ़ोरस पहले
समाप्त होता है और वृद्धि को सीमित करता है।
\end{solution}

\begin{solution}{exo:b3:environmental-toxicology:5}
$L_0 = \mathrm{BOD_5}/(1 - \eu^{-0.23 \times 5}) = 170/0.683 = \qty{250}{mg/L}$।
\end{solution}

\begin{solution}{exo:b3:environmental-toxicology:6}
\qty{0.1000}{L} में $\qty{4.40e-3}{L} \times \qty{0.0200}{mol/L} = \qty{8.80e-5}{mol}$: \omterm{def:b3:environmental-toxicology:alkalinity}{क्षारीयता} $\qty{8.80e-4}{mol/L}$; $\ce{HCO3-}$ (\qty{61.0}{g/mol}) के रूप में,
\qty{53.7}{mg/L}।
\end{solution}

\begin{solution}{exo:b3:environmental-toxicology:7}
$K_d = 0.020 \times 300 = \qty{6.0}{L/kg}$। शोषित/कुल $= K_dm_s/(K_dm_s + V_w) = 9.0/(9.0 + 0.30) = 0.97$: 97\,\% शोषित।
\end{solution}

\begin{solution}{exo:b3:environmental-toxicology:8}
बेंज़ीन: $\log\mathrm{BCF} = 0.85 \times 2.13 - 0.70 = 1.11$, BCF लगभग 13। PCB: $0.85 \times 6.67 - 0.70 = 4.97$, BCF लगभग
\num{9e4}: यह मछली में बेंज़ीन की तुलना में लगभग सात हज़ार गुना अधिक सांद्रित होता है; ऐसे
$K_{ow}$ पर समाश्रयण अपनी वैधता की सीमा के निकट है।
\end{solution}

\begin{solution}{exo:b3:environmental-toxicology:9}
$\qty{192}{mg/kg} \times \qty{70}{kg} = \qty{13}{g}$, लगभग 150 प्याले। प्रजातियाँ अवशोषण और उपापचय में भिन्न होती हैं,
और $\mathrm{LD_{50}}$ कोई सुरक्षित देहली नहीं है; हानिकारक प्रभाव उससे बहुत नीचे आरंभ हो जाते हैं।
\end{solution}

\begin{solution}{exo:b3:environmental-toxicology:10}
PNEC $= \qty{0.50}{mg/L}/1000 = \qty{0.50}{\micro g/L}$; $\mathrm{RQ} = 0.20/0.50 = 0.40$, 1 से कम: इस अनावरण पर कोई चिंता नहीं।
\end{solution}

\begin{solution}{exo:b3:environmental-toxicology:11}
प्लवक/जल: $\qty{0.015}{mg/kg}/\qty{2.0e-6}{mg/L} = \qty{7.5e3}{L/kg}$ (जैवसांद्रण)। मछली/प्लवक: $0.30/0.015 = 20$;
पक्षी/मछली: $6.0/0.30 = 20$ (प्रत्येक चरण पर \omterm{def:b3:environmental-toxicology:bio}{जैवआवर्धन}); पक्षी/जल: $\num{3e6}$।
\end{solution}

\begin{solution}{exo:b3:environmental-toxicology:12}
$\qty{1.0e9}{g}/\qty{64.1}{g/mol} = \qty{1.56e7}{mol}$ $\ce{SO2}$, जो उतने ही मोल $\ce{H2SO4}$ देता है: $\qty{1.53e9}{g}$, लगभग \qty{1500}{t}। यह
$\qty{3.12e7}{mol}$ $\ce{H+}$ मुक्त करता है; pH 4.0 पर $[\ce{H+}] = \qty{1.0e-4}{mol/L}$: $\qty{3.1e11}{L}$ वर्षा।
\end{solution}

\begin{solution}{pb:b3:environmental-toxicology:1}
\textbf{1.} यह जल की तुलना में ऑक्टेनॉल में दस हज़ार गुना अधिक विलेय है: जलविरागी होने से
यह कार्बनिक पदार्थ पर शोषित होगा और वसा में संचित होगा।

\textbf{2.} $K_{oc} = 0.41 \times 10^4 = \qty{4100}{L/kg}$; $K_{sw} = f_{oc}K_{oc} = 0.040 \times 4100 = \qty{164}{L/kg}$।

\textbf{3.} $K_{fw} = 0.048 \times 10^4 = \qty{480}{L/kg}$।

\textbf{4.} नहीं: साम्य पर वायु में जल-सांद्रता का केवल $10^{-4}$ भाग होता है।

\textbf{5.} उदासीन, जलविरागी अणु कार्बनिक पदार्थ में घुल जाता है, जैसे
ऑक्टेनॉल में; खनिज पृष्ठ ध्रुवीय होते हैं और जल से ढके रहते हैं।

\textbf{6.} तलछट, जिसकी क्षमता, $K_{sw}m_s$, सबसे बड़ी है।

\textbf{7.} $V_w = \qty{1.0e9}{L}$; $K_{aw}V_a = \qty{1.0e8}{L}$; $K_{sw}m_s = \qty{1.64e10}{L}$; $K_{fw}m_f = \qty{4.8e6}{L}$।

\textbf{8.} $C_w = \qty{1.0e8}{mg}/\qty{1.75e10}{L} = \qty{5.7e-3}{mg/L}$, \qty{5.7}{\micro g/L}।

\textbf{9.} जल \qty{5.7}{kg}, वायु \qty{0.57}{kg}, तलछट \qty{94}{kg}, मछली \qty{0.027}{kg}।

\textbf{10.} वायु $\qty{5.7e-7}{mg/L}$; तलछट \qty{0.94}{mg/kg}; मछली \qty{2.7}{mg/kg}।

\textbf{11.} सभी विभागों के बीच एक साथ साम्य, कोई निम्नीकरण नहीं, कोई
अंतर्वाह या बहिर्वाह नहीं, और भली-भाँति मिश्रित विभाग।

\textbf{12.} स्थायी अवस्था में निम्नीकरण और प्रवाहों के साथ (स्तर II) कुल भंडार
निवेश और हानि की दरों से तय होता है; विभागों के बीच स्थानांतरण दरों के साथ (स्तर III)
सांद्रताएँ अब साम्य अनुपातों में नहीं रहतीं, और निवेश प्राप्त करने वाला
विभाग समृद्ध हो जाता है।

\textbf{13.} $k = \ln 2/\qty{60}{d} = \qty{0.0116}{d^{-1}}$।

\textbf{14.} $\eu^{-0.0116 \times 365} = 0.0145$: 1.5\,\% शेष।

\textbf{15.} $\qty{5.7}{\micro g/L} \times 0.0145 = \qty{0.083}{\micro g/L}$।

\textbf{16.} ठंडे, ऑक्सीजनरहित तलछट में निम्नीकरण प्रायः धीमा होता है, और शोषित
पीड़कनाशी धीरे-धीरे वापस जल में मुक्त होता है, जिससे जल के साफ़ हो जाने के बाद भी
वह उपस्थित बना रहता है।

\textbf{17.} जल-अपघटन, ऑक्सीकरण और सूक्ष्मजीवी आक्रमण का प्रतिरोध करने वाले आबंध (उदाहरण के लिए
ऐरोमैटिक C--Cl), कम जल-विलेयता, और ऐसा शोषण जो उसे
निम्नीकारकों से बचाता है।

\textbf{18.} PNEC $= \qty{50}{\micro g/L}/10 = \qty{5.0}{\micro g/L}$।

\textbf{19.} $\mathrm{RQ} = 5.7/5.0 = 1.1$।

\textbf{20.} 1 से ठीक ऊपर: एक जोखिम संकेतित है, और मूल्यांकन को परिष्कृत करना होगा (बेहतर
विषाक्तता आँकड़े, मापी गई सांद्रताएँ) या अनावरण घटाना होगा।

\textbf{21.} \qty{2.7}{mg/kg}, \qty{1.0}{mg/kg} की देहली का लगभग तीन गुना: जब तक सांद्रता न गिरे,
मछली नहीं खाई जानी चाहिए।

\textbf{22.} समय के साथ और कई बिंदुओं पर जल, तलछट और मछली में सांद्रताएँ,
मॉडल की जाँच करने और गिरावट का अनुसरण करने के लिए।

\textbf{23.} कम पीड़कनाशी का प्रयोग, वर्षा से और तट से दूर प्रयोग,
अपवाह रोकने वाली बफ़र पट्टियाँ, या कम चिरस्थायी विकल्प।

\textbf{24.} यह मापे गए निष्प्रभाव स्तर को भाग देकर कुछ प्रयोगशाला-प्रजातियों और
संपूर्ण पारितंत्र के बीच के अंतर को समाहित करता है; 1000 के बजाय 10 का गुणक
PNEC को, और निर्णय को, सौ गुना बदल देता है।

\textbf{25.} \omterm{def:b3:environmental-toxicology:ecotoxicology}{जोखिम भागफल} PEC/PNEC लगभग 1.1 है, 1 से ठीक ऊपर।
\end{solution}
